\ECchapter{04}{Nuclear Energy}{Can nuclear power support a low-carbon future?}{ECBlue}

\ECObjectives{
\begin{itemize}[leftmargin=*,itemsep=1pt,topsep=1pt]
\item Explain nuclear fission and a controlled chain reaction.
\item Identify the energy and fluid flows in a pressurised water reactor (PWR).
\item Explain why residual heat must be removed after shutdown.
\item Interpret lifecycle-emission data and discuss the role of nuclear power.
\end{itemize}}

\begin{multicols}{2}

\begin{engineeringnews}
Several countries are considering new reactors, lifetime extensions and Small Modular Reactors (SMRs). The main engineering questions concern construction time, financing, safety, fuel supply, waste management and operation within an electricity system containing more variable renewable generation.
\end{engineeringnews}

\begin{ecarticle}
A nuclear power station converts energy released by \textbf{nuclear fission} into electricity. When a uranium-235 nucleus absorbs a neutron, it may split into two lighter nuclei. This reaction releases energy and additional neutrons, which can trigger further fissions. In a reactor, the number of fissions is controlled so that power remains stable.

In a pressurised water reactor, fuel assemblies are placed inside the reactor vessel. Water in the primary circuit flows through the core and removes heat. It is kept at high pressure so that it remains liquid at about 300~$^{\circ}$C. In the steam generator, the primary water transfers heat to a physically separate secondary circuit.

The secondary water boils and produces steam. The steam drives a turbine connected to an electrical generator. After the turbine, a condenser turns the steam back into liquid water. A third circuit carries waste heat towards a river, the sea or a cooling tower. The three circuits are separated because they have different functions and operating conditions.

Reactor power is adjusted by controlling neutron absorption. Control rods contain materials that absorb neutrons, while soluble boron can also be added to the primary water. Even after the chain reaction stops, radioactive decay continues to release \textbf{residual heat}. Cooling must therefore remain available after shutdown.

Safety is based on defence in depth: prevention, monitoring, automatic protection, backup systems and several physical barriers. Engineers also study component ageing, extreme weather, human factors, cybersecurity and long-term waste management.
\end{ecarticle}

\columnbreak

\begin{engineeringfocus}
\textbf{Energy conversion in a PWR}
\begin{center}
\begin{tikzpicture}[scale=.54,transform shape,>=Latex,font=\scriptsize]
\tikzset{block/.style={draw=ECBlue,thick,rounded corners,fill=ECLightBlue,minimum width=21mm,minimum height=10mm,align=center}}
\node[block] (core) {reactor\\core};
\node[block,right=9mm of core] (sg) {steam\\generator};
\node[block,right=9mm of sg] (turbine) {turbine};
\node[block,below=10mm of turbine] (cond) {condenser};
\node[block,below=10mm of core] (cool) {cooling\\source};
\draw[->,very thick,ECOrange] (core) -- node[above]{primary water} (sg);
\draw[->,very thick,ECOrange] (sg) to[bend left=22] node[below]{return} (core);
\draw[->,very thick,ECGreen] (sg) -- node[above]{steam} (turbine);
\draw[->,very thick,ECGreen] (turbine) -- (cond);
\draw[->,very thick,ECGreen] (cond) -| node[pos=.25,below]{feedwater} (sg);
\draw[->,very thick,ECBlue] (cool) -| (cond);
\draw[->,very thick,ECBlue] (cond) -| node[pos=.25,right]{warm water} (cool);
\draw[->,very thick] (turbine) -- ++(15mm,0) node[right]{generator};
\end{tikzpicture}
\end{center}
The primary circuit transports nuclear heat, the secondary circuit produces steam, and the cooling circuit rejects unused heat. A PWR is therefore both a nuclear system and a thermal power plant.
\end{engineeringfocus}

\begin{tcolorbox}[title=\sffamily\bfseries Did You Know?,colback=yellow!10,colframe=ECOrange]
Stopping fission does not stop heat production immediately. Residual heat decreases rapidly, but cooling remains essential for hours, days and longer after shutdown.
\end{tcolorbox}

\begin{tcolorbox}[title=\sffamily\bfseries Key Figures,colback=ECLightBlue,colframe=ECBlue]
\begin{tabularx}{\linewidth}{@{}Xr@{}}
Typical French PWR output & 900--1450 MW\\
Primary-water temperature & about 300~$^{\circ}$C\\
Primary pressure & about 155 bar\\
Typical capacity factor & 70--90\%\\
Operating life & several decades
\end{tabularx}
\end{tcolorbox}

\begin{tcolorbox}[title=\sffamily\bfseries Indicative lifecycle emissions,colback=white,colframe=ECBlue]
\centering
\begin{tikzpicture}
\begin{axis}[width=\linewidth,height=48mm,ybar,bar width=7pt,ymin=0,ymax=900,ylabel={g CO$_2$-eq/kWh},symbolic x coords={Nuclear,Wind,Solar PV,Natural gas,Coal},xtick=data,x tick label style={rotate=35,anchor=east,font=\scriptsize},yticklabel style={font=\scriptsize},ylabel style={font=\scriptsize},grid=major]
\addplot+[fill=ECBlue!45,draw=ECBlue] table[x=source,y=emissions,col sep=comma]{data/EC04/lifecycle_emissions.csv};
\end{axis}
\end{tikzpicture}
{\scriptsize Indicative median lifecycle values based on IPCC assessments. Actual values vary with technology, location and methodology.}
\end{tcolorbox}

\begin{tcolorbox}[title=\sffamily\bfseries Engineering Insight,colback=white,colframe=ECOrange]
A power station must remove more thermal energy than it converts into electricity. This is why condenser performance and access to a cooling source strongly influence plant design, site selection and output during heat waves.
\end{tcolorbox}

\begin{tcolorbox}[title=\sffamily\bfseries Fission and fusion,colback=white,colframe=ECBlue]
\textbf{Fission} splits heavy nuclei and is used in current commercial reactors. \textbf{Fusion} combines light nuclei and powers the Sun. Fusion research is active, but fusion is not yet a commercial source of electricity.
\end{tcolorbox}

\begin{vocabulary}
\begin{tabularx}{\linewidth}{@{}lX@{}}
chain reaction & réaction en chaîne\\
reactor vessel & cuve du réacteur\\
fuel assembly & assemblage combustible\\
control rod & barre de contrôle\\
steam generator & générateur de vapeur\\
containment & enceinte de confinement\\
residual heat & puissance résiduelle\\
spent fuel & combustible usé
\end{tabularx}
\end{vocabulary}

\begin{questions}
\item What happens during nuclear fission?
\item Why must the primary water remain under high pressure?
\item Explain the function of each of the three water circuits.
\item How do control rods modify reactor power?
\item Why is cooling required after shutdown?
\item Use the graph to compare nuclear power with fossil-fuel generation.
\item Why can high air or river temperatures reduce the output of a power station?
\end{questions}

\begin{challenge}
A region must replace 8~GW of fossil-fuelled generation while reducing lifecycle CO$_2$ emissions and maintaining supply during winter evenings. Propose a balanced mix using nuclear, wind, solar, storage, interconnections and demand response. Justify installed capacity, expected availability, flexibility, environmental constraints, construction sequence and major risks.
\end{challenge}

\begin{applicationexercise}
A 1.3~GW nuclear unit operates with a capacity factor of 88\%. Estimate its annual electrical production in TWh. Then determine how many 3~MW wind turbines operating at a 32\% capacity factor would produce the same annual energy. Explain why equal annual production does not mean that the two systems provide electricity in the same way.
\end{applicationexercise}

\begin{speaking}
Debate: ``Europe should build new nuclear power stations.'' Use evidence about safety, climate, cost, waste, construction time and security of supply. Each team must answer one technical objection from the opposing side.
\end{speaking}

\end{multicols}

\subsection*{Sources and further reading}
\ECsource{IAEA -- Nuclear Power Reactors}{https://www.iaea.org/topics/nuclear-power-reactors}\quad
\ECsource{OECD Nuclear Energy Agency}{https://www.oecd-nea.org}\quad
\ECsource{IPCC -- Energy Systems}{https://www.ipcc.ch}
